\documentclass[11pt]{article}
\usepackage[margin=1in]{geometry}
\usepackage{amsmath,amssymb,amsthm,mathtools,microtype}
\usepackage[hidelinks]{hyperref}

\newtheorem{theorem}{Theorem}[section]
\newtheorem{lemma}[theorem]{Lemma}
\newtheorem{proposition}[theorem]{Proposition}
\newtheorem{corollary}[theorem]{Corollary}
\newtheorem{remark}[theorem]{Remark}
\newcommand{\A}{\mathbb A}
\renewcommand{\P}{\mathbb P}
\newcommand{\C}{\mathbb C}

\title{Decorated Discriminant Normalization and Stable Moduli}
\author{Repository working draft}
\date{\today}

\def\DecoratedNormalizationPaper{1}

\begin{document}
\maketitle

\begin{abstract}
For the weighted marked-root Keller maps attached to a normalized seed $H$,
we study the normalization of the discriminant curve of the pencil
$H(W)-sW+t$.  Its zeroth relative Fitting ideal is $(H''(r))$ with full
multiplicity.  We retain the saturated off-diagonal node-pairing scheme, the
conductor morphism, the point at infinity, and the scheme-valued boundary
marks.  On a nonempty ordinary boundary-clean open, the resulting stable
invariant is generically etale of degree $N-2$ onto a normalized image of
dimension $N-3$.  Consequently degree-$N$ weighted maps contain an
$(N-3)$-dimensional family of stable classes.  After adding the unique
unramified affine root sheet, the normalized seed open is the normalization
of the marked decorated image; forgetting that sheet has degree $N-2$.
\end{abstract}

\section{Setup and scope}

Let $k$ be a characteristic-zero field and let $H\in k[W]$ be a normalized
admissible degree-$N$ seed:
\[
 H(0)=H'(0)=H(1)=0,\qquad H'(1)=-1.
\]
The weighted construction associates to $H$ a polynomial Keller map $F_H$
and the marked inverse equation
\[
 E_H(W;s,t)=H(W)-sW+t=0.
\]
Its finite incidence normalization is equipped with the maximal open on which
the rational source reconstruction is regular.  The construction-independent
\href{../../verified/STABLE_NORMALIZATION_FUNCTORIALITY.md}
{stable-normalization functoriality theorem} identifies the
intrinsic discriminant vertex, its finite normalization morphism, and its
boundary marks under polynomial left--right equivalence and after adjoining
identity variables.  Those facts are the functorial input used here.

The coarse decorated invariant deliberately forgets the affine root sheet,
so its generic degree is the rerooting degree $N-2$.  Retaining the unique
unramified sheet meeting the regular-reconstruction open removes exactly this
ambiguity and gives generic faithfulness.  Extension of that marking through
arbitrary collision strata remains a separate valuative problem; none of the
Fitting, node, conductor, generic-degree, or stable-moduli results depends on
that extension.

\section{The scheme-theoretic decorated normalization}

For a weighted primitive $H$, let $D_H$ be the reduced discriminant of
$H(W)-sW+t$ and let
\[
 \nu_H:\A^1_r\longrightarrow D_H,
 \qquad r\longmapsto(H'(r),rH'(r)-H(r))                 \tag{6.1}
\]
be its affine normalization.  Write
$R_H=k[H'(r),rH'(r)-H(r)]\subset k[r]$.

\begin{proposition}[full Fitting divisor]
\label{prop:full-fitting}
There is an equality of ideal sheaves on the normalization
\[
 \operatorname{Fitt}_0\Omega_{\widetilde D_H/D_H}=(H''(r)). \tag{6.2}
\]
This ideal, including every root multiplicity, is preserved by polynomial
left--right equivalence and by adjoining identity variables.
\end{proposition}

\begin{proof}
The images in $k[r]\,dr$ of the two generators of $R_H$ are
\[
 dH'(r)=H''(r)dr,
 \qquad d(rH'(r)-H(r))=rH''(r)dr.
\]
Thus
$\Omega_{k[r]/R_H}\simeq k[r]/(H'')\,dr$, proving (6.2).
The stable boundary theorem identifies the finite normalization morphisms on
the intrinsic discriminant vertex away from the second vertex, so it
identifies their relative differential modules and Fitting ideals.  Under
stabilization one has
\[
 \Omega_{(\widetilde D_H\times\A^q)/(D_H\times\A^q)}
 \simeq\operatorname{pr}_1^*\Omega_{\widetilde D_H/D_H}.
\]
Zeroth Fitting ideals commute with this flat base change.  Equality of the
ideal sheaves preserves their valuations at every height-one point, hence
preserves the multiplicities in (6.2), not only the reduced support.
\end{proof}

If the support of (6.2) contains two points, the units of
$k[r,\xi,\xi^{-1},z_1,\ldots,z_q]$ and subtraction of the two corresponding
principal-prime equations force the induced coordinate change to have the
form $r\mapsto ar+b$, $a\ne0$.  Proposition~\ref{prop:full-fitting} then says
that the full effective Hessian divisor is carried affinely, with its
coefficients.

The self-intersection of (6.1) is also intrinsic.  In two normalization
coordinates put
\[
 D_s(r,u)=\frac{H'(r)-H'(u)}{r-u},
\]
\[
 D_t(r,u)=
 \frac{rH'(r)-H(r)-uH'(u)+H(u)}{r-u}.
\]
The genuine off-diagonal closure is the scheme
\[
 \Gamma_H^{\rm off}
 =V\bigl((D_s,D_t):(r-u)^\infty\bigr).                 \tag{6.3}
\]
The saturation in (6.3) is necessary: on the diagonal the divided equations
specialize to $H''(r)$ and $rH''(r)$ and otherwise retain isolated cusp
points.  The involution $(r,u)\leftrightarrow(u,r)$ is free on the ordinary
off-diagonal locus; its characteristic-zero finite quotient exists in
general and retains fixed diagonal limits in degenerations.  On the ordinary
locus that quotient records the two normalization branches paired at each node.  An
isomorphism of normalization morphisms preserves the fiber product, its
diagonal, and the saturated complement, while stabilization takes their
product with affine space.  Hence this node-pairing scheme is stable.
In the exact extractor, ordinary-node status is certified by the Jacobian
determinant
\[
 \det\frac{\partial(D_s,D_t)}{\partial(r,u)}
\]
being a unit on (6.3), rather than inferred only from the number of projected
branch points.  Passing to
$\sigma=r+u,\pi=ru$ gives the scheme quotient by the free involution; the
implementation verifies that its equations pull back into the saturated
ordered ideal.

Finally, let
\[
 \mathfrak c_H=\operatorname{Ann}_{\mathcal O_{D_H}}
 (\nu_{H*}\mathcal O_{\widetilde D_H}/\mathcal O_{D_H}) \tag{6.4}
\]
be the conductor.  It is functorial for the same reason.  On the open locus
of ordinary cusps and nodes, if $P_H(r)$ cuts out all normalization branches
appearing in (6.3), then the local models
$k[[t^2,t^3]]\subset k[[t]]$ and
$k[[x,y]]/(xy)\subset k[[x]]\oplus k[[y]]$ give
\[
 \mathfrak c_Hk[r]=(H''(r)^2P_H(r)).                   \tag{6.5}
\]
For arbitrary plane-curve singularities, let $f_H(S,T)=0$ be the primitive
implicit equation.  Plane-curve adjunction gives the uniform generator
\[
 \mathfrak c_Hk[r]=
 \left(
 \frac{(\partial f_H/\partial T)
 (H'(r),rH'(r)-H(r))}{H''(r)}
 \right).                                             \tag{6.5a}
\]
This quotient is polynomial because (6.1) is the normalization.  The exact
extractor uses (6.5a) without an ordinary-singularity hypothesis and checks
(6.5) independently whenever the cusp/node criterion holds.
The invariant retains the finite conductor morphism
\[
 \operatorname{Spec}(k[r]/\mathfrak c_Hk[r])
 \longrightarrow
 \operatorname{Spec}(R_H/(R_H\cap\mathfrak c_Hk[r])), \tag{6.6}
\]
not only its upstairs support.  It records the two branches over a node and
the single thick branch over a cusp.  Together, (6.2)--(6.6), the point at infinity, and the labeled inverse image
of the intersection with the second target boundary form the decorated
normalization.  The full second-boundary center divisor upstairs is
$\gcd(H,H')$; a point label is inferred only when its support is uniquely
certified.  Every part is therefore preserved under stable polynomial
left--right equivalence.

For the split quartics, the boundary label on the normalization is also
checked intrinsically.  If zero is the unique multiple primitive root and
$H(r)=h_2r^2+O(r^3)$, the normalized discriminant boundary chart
\[
 B=\frac{H'(r)}C,\qquad
 A=\frac{rH'(r)-H(r)}{cC^2},\qquad r=Ck
\]
specializes to
\[
 B=2h_2k,\qquad A=\frac{h_2}{c}k^2.
\]
A finite specialization centered at $r=\rho$ forces
$H(\rho)=H'(\rho)=0$, hence $\rho=0$.  Thus the second-boundary intersection
canonically marks $r=0$; it is not an auxiliary normalization convention.

For completeness, this decoration gives an actual generic-dimension theorem,
not only a parameter count.  In the normalized degree-$N$ seed space, an
isomorphism preserving zero and infinity is $r\mapsto ar$.  Equality of
Fitting divisors then has the form $G''(r)=cH''(ar)$.  Since the normalized
seeds vanish together with their first derivatives at zero, integration gives
\[
 G(r)=\frac{c}{a^2}H(ar).
\]
The condition $G(1)=0$ forces $H(a)=0$, leaving at most $N-2$ generic
choices, and $G'(1)=-1$ determines $c$.  Thus the decorated-normalization
moduli map is generically finite and its image has dimension $N-3$.


\end{document}
